% Auto-generado por scripts/extract-lessons-to-tex.ps1
% Fuente: HTML en carpeta L4-Zonas-grises-y-ambiguedades-interpretativas (sin modificar archivos originales).

\documentclass[11pt,a4paper]{article}
\usepackage[utf8]{inputenc}
\usepackage[T1]{fontenc}
\usepackage[spanish]{babel}
\usepackage{geometry}
\geometry{margin=2.5cm}
\usepackage{hyperref}
\hypersetup{colorlinks=true, linkcolor=blue, urlcolor=blue}

\title{\texttt{L4-Zonas-grises-y-ambiguedades-interpretativas}}
\author{Extracci\'on desde lecciones HTML}
\date{}

\begin{document}
\maketitle

\begin{small}
\noindent Este documento consolida el texto visible de los ocho componentes (01--08). Los formularios interactivos aparecen como texto; la l\'ogica JavaScript no se reproduce aqu\'i.
\end{small}
\bigskip
\section{Conceptos}
\noindent\textit{Archivo fuente:} \texttt{01-conceptos.html}

L4 — Conceptos\par





    L4 — Elemento 1 — Conceptos\par
    Tres zonas grises del temario\par

    Introducción breve: el temario identifica áreas que exigen posición institucional explícita en el curso.\par


      Esta lección sistematiza las tres fricciones descritas en el PDF del temario para revisión del experto (SME).\par
      No sustituye la decisión académica de su programa: ofrece mapa conceptual y consecuencias procesales (preclusión, archivo,\par
      casación) para que el litigante elija línea con plena conciencia del riesgo.\par



      1. Falta de lesividad / insignificancia (bagatela)\par


        Conflicto: parte de la jurisprudencia de la CSJ ha tratado la ausencia de daño relevante al bien jurídico como\par
        cuestión de atipicidad objetiva (con efectos en causal de preclusión distinta), mientras otra línea la ubica en el\par
        plano de la antijuridicidad material. El temario cita, entre otras, Rad. 21064, 31362 frente a 47862, 51816.\par


      Efecto práctico: según la línea elegida, variará la causal de preclusión invocable (p. ej. art. 332 CPP: causal 4 atipicidad vs. causal 6 antijuridicidad) y la viabilidad del archivo fiscal.\par




      2. Pena natural (poena naturalis) — art. 34 C.P.\par


        Figura por la cual el daño gravísimo sufrido por el autor puede conducir a exención de pena cuando concurran los\par
        requisitos legales. El temario señala desarrollo jurisprudencial mayoritario en culpa (p. ej. accidentes con resultado\par
        en familiares) y debate sobre extensión analógica in bonam partem a supuestos dolosos.\par





      3. Doble subsunción y crímenes internacionales\par


        Tensión entre subsumir el mismo hecho en tipo interno y en figuras del Estatuto de Roma (p. ej. lesa humanidad),\par
        frente a críticas de legalidad estricta, lex praevia y analogía prohibida que agrave (doctrina citada en temario:\par
        Malarino, Gálvez; caso César Pérez García).\par





      4. Cuadro resumen\par



          Zona grisNúcleo problemáticoDecisión del curso (parametrizable)\par




          LesividadAtipicidad vs. antijuridicidad materialDefinir línea pedagógica y modelo de escrito\par

          Pena naturalAlcance en doloDelimitar tesis defensiva prudente\par

          Doble subsunciónCP + art. 7 ERElegir postura garantista u homogeneizadora con CSJ\par








      Fuentes: Temario Propuesta R (1).pdf, sección «Zonas grises y ambigüedades interpretativas». Bernal (2024), Cap. 4 y material del temario.\par


    Siguiente →\par

\section{Explicaci\'on}
\noindent\textit{Archivo fuente:} \texttt{02-explicacion.html}

L4 — Explicación\par





    L4 — Elemento 2 — Explicación\par
    Explicación dogmática y procesal\par


      Cada zona gris exige doble lectura: (i) ubicación en la teoría del delito; (ii) traducción a actos procesales concretos\par
      (solicitud de preclusión, recurso de casación, acción de inconstitucionalidad).\par



      1. Lesividad: dos mapas posibles\par



          Línea “atipicidad objetiva”\par

          La insignificancia vacía el contenido típico del delito en el caso concreto; encaja en lógica de causal 4 (art. 332 CPP) según el temario.\par




          Línea “antijuridicidad material”\par

          El tipo se cumple formalmente pero falta lesión o menoscabo relevante; encaje procesal asociado a causal 6 y límites al archivo fiscal según el temario.\par








      2. Art. 34 C.P. y test de necesidad\par


        La exención por pena natural se articula con principio de necesidad y proporcionalidad (arts. 3 y 4 C.P. en el discurso del temario).\par
        La defensa debe demostrar el daño gravísimo al autor y el desplazamiento de la necesidad de la pena estatal\par
        frente a la “pena vivida”. La extensión a dolo es argumento de alto riesgo que requiere anclaje jurisprudencial fino.\par





      3. Doble subsunción y legalidad\par


        La tesis crítica sostiene que la doble calificación puede vulnerar lex praevia y prohibiciones de analogía in malam partem,\par
        o generar concursos atípicos con efectos acumulativos inconstitucionales. La tesis defensiva deberá separar\par
        competencia, tipicidad y ne bis in idem material antes de entrar al fondo sancionatorio.\par





      Fuentes: Temario (anexo zonas grises); Guzmán Díaz (2021) sobre ausencia de lesividad; Boada Acosta (2019) pena natural.\par


    Siguiente →\par

\section{Exploraci\'on}
\noindent\textit{Archivo fuente:} \texttt{03-exploracion.html}

L4 — Exploración\par





    L4 — Elemento 3 — Exploración\par
    Elegir línea procesal bajo incertidumbre\par

    Tres microcasos; la “mejor” respuesta depende de la línea institucional que adopte su curso o su cliente.\par


      Caso A — Hurto de mínima cuantía\par

      Defensa pide cierre temprano por insignificancia. ¿Invoca primero atipicidad (4) o antijuridicidad (6)?\par

       Solo causal 4\par
       Solo causal 6\par
       Estrategia dual con subsidiariedad en el escrito\par
      Retroalimentación\par





      Caso B — Dolo con consecuencia catastrófica para el autor\par

      ¿Solicitaría art. 34 en delito doloso?\par

       Sí, por analogía in bonam partem\par
       No, reservado a culpa según práctica dominante\par
       Subsidiariamente, con dictamen pericial del daño al autor\par
      Retroalimentación\par





      Caso C — Imputación simultánea CP + lesa humanidad\par

       Aceptar doble subsunción si la CSJ la validó\par
       Atacar por legalidad y ne bis in idem material\par
       Negociar desistimiento de una vía\par
      Retroalimentación\par




    Siguiente →\par

\section{Contexto}
\noindent\textit{Archivo fuente:} \texttt{04-contexto.html}

L4 — Contexto\par





    L4 — Elemento 4 — Contexto\par
    CSJ, Fiscalía y defensa en escenarios friccionados\par


      La fragmentación jurisprudencial no es un “detalle académico”: condiciona expectativas de archivo, ritmo procesal\par
      y estrategia de segunda instancia. El temario pide al experto fijar la línea que el curso enseñará como regla principal.\par





          ActorInterés típicoRiesgo si no se clarifica la dogmática\par




          FiscalíaPersecución efectivaArchivo indebido o denegada preclusión según causal mal escogida\par

          DefensaCierre o absoluciónPérdida de oportunidad por mala cualificación dogmática\par

          CSJ / casaciónCorrección de errores jurídicosLitigio repetido sobre mismos debates (21064 vs 47862)\par








      Justicia transicional y ne bis in idem\par


        El temario enlaza el debate de doble subsunción con foros internacionales y la Ley 906 (art. 192 citado en L1). La defensa debe mapear\par
        identidad fáctica y identidad jurídica antes de aceptar doble enjuiciamiento.\par




    Fuentes: Temario; CSJ Rad. citadas en el PDF.\par

    Siguiente →\par

\section{Aplicaci\'on}
\noindent\textit{Archivo fuente:} \texttt{05-aplicacion.html}

L4 — Aplicación\par





    L4 — Elemento 5 — Aplicación\par
    Clínica: tres minutas tipo\par

    Esquemas para adaptar al expediente; no sustituyen asesoría en caso concreto.\par


      Minuta 1 — Preclusión por insignificancia\par



-- Relación sintética del hecho y del tipo imputado.\par


-- Subsección A (subsidiaria): tesis de atipicidad objetiva — cita a línea 21064/31362.\par


-- Subsección B (subsidiaria): tesis de falta de antijuridicidad material — cita a línea 47862/51816.\par


-- Petitorio: preclusión causal 4 y/o 6 según subsección; archivo si procede.\par





      Minuta 2 — Solicitud de prescindencia (art. 34 C.P.)\par



-- Descripción del daño gravísimo al autor (prueba médica, psicológica, familiar).\par


-- Conexión con arts. 3 y 4 C.P. (necesidad y proporcionalidad de la pena estatal).\par


-- Distinción culpa/dolo y jurisprudencia de tutela (Rad. 48820 y similares).\par


-- Petitorio: absolución o exención de pena; subsidiarias.\par





      Minuta 3 — Casación / apelación por doble subsunción\par



-- Identidad del hecho juzgado en CP y en figura internacional.\par


-- Agravio: vulneración de lex praevia, tipicidad estricta o ne bis in idem.\par


-- Cita doctrinal garantista (Malarino, Gálvez) como corolario, no sustituto de prueba.\par


-- Petitorio: anular parcialmente / unificar delito / revocatoria de pena.\par




    Fuentes: Temario; CPP art. 332; C.P. arts. 3, 4, 34.\par

    Siguiente →\par

\section{Prospectiva}
\noindent\textit{Archivo fuente:} \texttt{06-prospectiva.html}

L4 — Prospectiva\par





    L4 — Elemento 6 — Prospectiva\par
    Tendencias y unificación\par


      Es plausible que la CSJ o la Corte Constitucional avancen hacia mayor unificación sobre lesividad y sobre límites de la\par
      cooperación penal internacional. El litigante debe monitorear SU-167/2024 y boletines de casación citados en bibliografía del temario.\par





          VariableEvolución posiblePreparación\par




          LesividadRegla única causal 4 o 6Actualizar plantillas de preclusión\par

          Art. 34Clarificación dolo/culpaBanco de precedentes filtrados por ratio\par

          MacrocrimenMás usos del Estatuto de RomaEquipo transnacional y amicus\par







    Fuentes: Temario bibliografía; Corte Constitucional SU-167/2024.\par

    Siguiente →\par

\section{Ejercicios}
\noindent\textit{Archivo fuente:} \texttt{07-ejercicios.html}

L4 — Ejercicios\par





    L4 — Elemento 7 — Ejercicios\par
    Ejercicios\par


      1. Redacte una oración que explique por qué la elección entre causal 4 y 6 no es indiferente para la Fiscalía.\par

      Modelo\par





      2. Enumere tres pruebas que apoyen art. 34 C.P.\par

      Modelo\par





      3. ¿Qué argumento de legalidad estricta puede usarse contra doble subsunción CP + art. 7 ER?\par

      Modelo\par





      4. Diferencie ne bis in idem interno e internacional en una línea.\par

      Modelo\par





      5. Mencione dos sentencias CSJ citadas en el temario sobre lesividad.\par

      Modelo\par





      6. ¿Qué papel tiene Guzmán Díaz (2021) según bibliografía del temario?\par

      Modelo\par




    Fuentes: Temario Propuesta R.\par

    Quiz →\par

\section{Quiz}
\noindent\textit{Archivo fuente:} \texttt{08-quiz.html}

L4 — Quiz\par





    L4 — Elemento 8 — Quiz\par
    Evaluación L4: zonas grises\par


      1. El temario señala fricción sobre “falta de lesividad” entre líneas que la ubican en:\par

         A) Solo derecho civil.\par
         B) Atipicidad objetiva frente a antijuridicidad material.\par
         C) Solo ejecución penal.\par


      2. Según el temario, el art. 332 CPP — causal 4 — se asocia típicamente a:\par

         A) Atipicidad.\par
         B) Solo error judicial.\par
         C) Solo prescripción.\par


      3. La causal 6 del art. 332 CPP, en el debate del temario, se vincula a:\par

         A) Falta de antijuridicidad material.\par
         B) Solo multas.\par
         C) Competencia militar.\par


      4. La poena naturalis en Colombia se regula principalmente en:\par

         A) Art. 34 C.P.\par
         B) Art. 100 C.P.\par
         C) Solo reglamento carcelario.\par


      5. El temario indica que la jurisprudencia de pena natural ha sido dominante en:\par

         A) Delitos culposos con daño gravísimo al autor.\par
         B) Solo delitos políticos.\par
         C) Contrabando aduanero leve.\par


      6. Extender art. 34 a dolo por analogía in bonam partem es según el temario:\par

         A) Tema debatido que requiere validación institucional.\par
         B) Imposible sin excepción.\par
         C) Resuelto sin debate.\par


      7. La doble subsunción CP + Estatuto de Roma tensiona principalmente:\par

         A) Legalidad estricta y límites interpretativos.\par
         B) Solo trámites de notificación.\par
         C) Solo honorarios.\par


      8. El caso citado en el temario vinculado a doble subsunción incluye:\par

         A) César Pérez García.\par
         B) Caso ficticio “Chavero”.\par
         C) Solo mediación civil.\par


      9. Doctrina citada en el temario como crítica garantista incluye:\par

         A) Malarino y Gálvez.\par
         B) Solo autores del siglo XIX.\par
         C) Solo normas tributarias.\par


      10. El anexo “Zonas grises” del temario está destinado a:\par

         A) Revisión del experto (SME) para fijar postura del curso.\par
         B) Solo diseño gráfico.\par
         C) Eliminar el Cap. 4 de Bernal.\par


      11. Boada Acosta (2019) en bibliografía del temario comenta:\par

         A) Pena natural y sentencia CSJ 2019.\par
         B) Solo derecho agrario.\par
         C) Solo arbitraje comercial.\par


      12. SU-167/2024 aparece en bibliografía como:\par

         A) Sentencia de unificación sobre garantías procesales.\par
         B) Código de tránsito.\par
         C) Ley de presupuesto.\par


      Enviar\par


      Resultado\par









    Inicio L4\par

\end{document}

